% Global: interpret evidence within the tested domain; local: no placeholder numerical claims.
\section{Discussion, claims, and implications}
The experiment separates a practical architecture comparison from a mechanistic question. A Standard--Block difference estimates the total effect of replacing residual composition under the matched training protocol. The Uniform Block control helps assess whether learned selection contributes beyond the change to normalized averaging, but its three-seed sample is exploratory. Interventions on trained Block weights examine immediate routing dependence while retaining the learned backbone; they are not interchangeable with training a uniform-routing model from scratch.

Until all primary pairs are evaluated, this manuscript establishes an implementation and a prospective test, not a forecast-quality improvement. A negative or inconclusive result remains informative about this particular low-budget, coarse-grid regime. Neither a positive local result nor a null result establishes behavior at foundation-model scale, unseen physical systems, different data volumes, or other residual block sizes.

\section{Limitations}
Six independent pairs constrain both power and the precision of variance estimates. Three-seed controls cannot support broad systematic claims. The confidence interval is conditional on selected trajectories and its normality approximation; no claim of coverage under arbitrary effect distributions is made. Matching optimizer updates does not match wall-clock compute because routing architectures have different costs. Conversely, equal compute would change the data exposure and is not the primary comparison specified here.

Coarsening removes unresolved spatial structure and changes numerical divergence residuals. Training-standardized NMSE can strongly weight low-energy targets despite the fixed stabilizer. The model learns a mixture of temporal strides without an explicit stride token, preserving the upstream forward interface; this may limit temporal identifiability. A 12-step endpoint is not a guarantee of asymptotic dynamical stability. GPU execution can vary across hardware and library versions even with deterministic algorithms requested\cite{pytorchrepro}. All these limitations apply regardless of the sign of the observed effect.

\section{Conclusion}
The study defines a reproducible test of depth-wise learned residual routing in the released GPhyT architecture, with matched seeds, controlled source scheduling, held-out autoregressive evaluation, and a finite compute budget. The accompanying measured-results section determines which accuracy and stability conclusions, if any, the completed runs support. Execution correctness and prospective power alone do not answer the primary scientific question.
